\documentclass[12pt]{article}
\usepackage[T1]{fontenc}
\usepackage[utf8]{inputenc}
\usepackage{lmodern}
\usepackage{textcomp}
\usepackage{enumitem}
\usepackage{geometry}
\geometry{margin=1in}
\begin{document}
\begin{center}
Answer Key - Physics - Graduate Level Assessment 13 \
\small (Answers are ordered exactly as in the question paper.)
\end{center}

\section*{Multiple Choice}
\begin{enumerate}[label=\arabic*.]
\item \textbf{Answer:} A
\\ \textit{Options:} A) 11.3; B) 7.0; C) 13.7; D) 6.5; E) 12.5
\\ \textit{Note:} The correct answer is 11.3 because the heat transfer by radiation is inversely proportional to the reflectance, and reducing the reflectance from 0.9 to 0.3 increases the effective emissivity, leading to approximately an 11.3-fold increase in heat transfer according to the Stefan-Boltzmann law.
\item \textbf{Answer:} A
\\ \textit{Options:} A) volume; B) timbre; C) speed; D) duration; E) wavelength
\\ \textit{Note:} The correct answer is volume because as the fire engine approaches, the sound waves compress due to the Doppler effect, resulting in an increase in amplitude and perceived loudness.
\item \textbf{Answer:} E
\\ \textit{Options:} A) 5.0 \ensuremath{\times} $10^3$ cal; B) 15.3 \ensuremath{\times} $10^3$ cal; C) 10.5 \ensuremath{\times} $10^3$ cal; D) 12.7 \ensuremath{\times} $10^3$ cal; E) 7.9 \ensuremath{\times} $10^3$ cal
\\ \textit{Note:} The correct answer of 7.9 × $10^3$ cal is derived from the calculation of electrical power using the formula P = VI, where the power (P) in watts is 550 watts (5.0 A × 110 V), which when converted to calories results in approximately 7.9 × $10^3$ cal for 1 minute, since 1 watt is equivalent to about 0.239 calories per second.
\item \textbf{Answer:} A
\\ \textit{Options:} A) 6.7 \ensuremath{\times} 10^-17 kg; B) 2.2 \ensuremath{\times} 10^-22 kg; C) 3.0 \ensuremath{\times} 10^-32 kg; D) 5.5 \ensuremath{\times} 10^-31 kg; E) 1.76 \ensuremath{\times} $10^11$ kg
\\ \textit{Note:} The correct answer of 6.7 × 10^-17 kg for the mass of the electron is derived by dividing Millikan's measured charge of the electron by Thomson's measured charge-to-mass ratio (q/m), accurately combining these fundamental constants to yield the electron's mass.
\item \textbf{Answer:} A
\\ \textit{Options:} A) remains the same; B) increases by 4; C) reduces by 3; D) increases by 3; E) increases by 1
\\ \textit{Note:} The correct answer is that the mass number of the resulting element remains the same because the ejection of an alpha particle, which consists of 2 protons and 2 neutrons, effectively results in the loss of 4 nucleons, but the element undergoing the decay is transformed into a different element with a new atomic number, maintaining the overall mass number in the context of nuclear reactions.
\end{enumerate}

\section*{True / False}
\begin{enumerate}[label=\arabic*.]
\setcounter{enumi}{5}
\item \textbf{Answer:} True
\\ \textit{Options:} A) True; B) False
\\ \textit{Note:} The statement is true because according to the work-energy theorem, the work done on a particle is equal to the change in its kinetic energy, and since force is defined as the rate of change of momentum, a constant force of 5.0 N will result in a corresponding rate of change in kinetic energy over time.
\item \textbf{Answer:} True
\\ \textit{Options:} A) True; B) False
\\ \textit{Note:} The statement is true because if a capacitor loses half of its charge every second, after five seconds it retains 1/32 of its initial charge, meaning if it has charge q after five seconds, its initial charge must have been 16 q (since q = (1/$2)^5$ * initial charge).
\item \textbf{Answer:} True
\\ \textit{Options:} A) True; B) False
\\ \textit{Note:} The statement is true because the magnification of an image in optics is defined as the ratio of the height of the image to the height of the object, and a magnification of 4 indicates that the image is indeed four times larger than the actual object.
\item \textbf{Answer:} False
\\ \textit{Options:} A) True; B) False
\\ \textit{Note:} The statement is false because, according to Newton's second law (F=ma), a force of 0.20 newton acting on a mass of 0.1 kg (100 grams) would result in an acceleration of 2.0 m/sec², not 5.0 m/sec².
\item \textbf{Answer:} False
\\ \textit{Options:} A) True; B) False
\\ \textit{Note:} Huygens' principle describes light as a wave phenomenon, illustrating how each point on a wavefront serves as a source of secondary wavelets, thereby demonstrating that light travels as waves rather than discrete particles.
\end{enumerate}

\section*{Long Form Response}
\begin{enumerate}[label=\arabic*.]
\setcounter{enumi}{10}
\item \textbf{Answer:} When a ballistic rocket is launched from the equator at a 60^\ensuremath{\circ} angle with an initial speed of 6 km/sec, it follows a parabolic trajectory influenced by gravitational forces and the curvature of the Earth. The great circle distance it covers can be calculated using the initial speed, launch angle, and the radius of the Earth. The horizontal component of the velocity can be determined using the cosine of the launch angle, resulting in a distance covered of approximately 8.660 \ensuremath{\times} $10^6$ m. This distance reflects the interplay between kinetic energy, gravitational potential energy, and the circular geometry of the Earth, demonstrating fundamental principles of projectile motion and orbital dynamics in a rotating reference frame.
\\ \textit{Note:} correct because it accurately applies the principles of projectile motion and gravitational forces to determine the horizontal distance traveled by the rocket, factoring in the launch angle and the Earth's curvature, which is essential for calculating the great circle distance.
\item \textbf{Answer:} In a head-on collision between a particle of mass \( m \) moving with velocity \( u_1 \) and a stationary particle of mass \( 2 m \), the dynamics can be analyzed using the principles of conservation of momentum and the coefficient of restitution. The conservation of momentum gives us the equation \( mu_1 = mv_1 + 2 mv_2 \), where \( v_1 \) and \( v_2 \) are the final velocities of the first and second particles, respectively. The coefficient of restitution \( e \) relates the relative velocities before and after the collision as \( e = \frac{v_2 - v_1}{u_1 - 0} \). To maximize the kinetic energy loss, we can set \( e \) to its lowest value, which is zero for a perfectly inelastic collision, leading to a significant change in velocity. Ultimately, for a perfectly elastic collision, we find that the final velocity \( v_1 \) of the first particle can be expressed as \( v_1 = \frac{1}{3}u_1 \), indicating that a higher coefficient of restitution results in greater retention of kinetic energy in the system.
\\ \textit{Note:} correct because it effectively applies the conservation of momentum and the definition of the coefficient of restitution to analyze the dynamics of the collision, demonstrating how these principles influence the final velocities and the kinetic energy loss during the interaction.
\end{enumerate}

\vfill
\noindent\textit{End of Answer Key}
\end{document}